% ===== 原文 PDF pp.93--97（印刷 pp.92--96）=====
% tr-08 / 块 tag 5.28--5.36 / 文件 c
% 起点：p.93 的 bullet 列表（p.92→p.93 的跨页句已由 tr08b 写完整）
% 终点：p.97 上部「致谢」段末；其后 References 由 Lead 脚本生成，本文件一字未写文献条目。
% 本文件含：表 8（p.94，秩一 SCFT 表）、§6 结论与开放问题、致谢。
% 说明：原文 §6 的标题实际位于 PDF p.93 底部（印刷 p.92，见页图），而表 8 作为浮动体排在 p.94 页首；
%       为保持「先见引用、后见表」的阅读顺序，表 8 按其在正文中的引用点（p.93 末「汇总在表 8 中」）就近排在 §6 标题之前。

% [原文 p.92]
\begin{itemize}
  \item $k=2$ 且 $\Delta_7=\frac{3}{2}, 2, 3$，相应的 7-膜为 $IV$、$I_0^*$ 与 $IV^*$ 型。
  \item $k=3$ 且 $\Delta_7=\frac{4}{3}, 2$，相应的 7-膜为 $III$、$I_0^*$ 型。
  \item $k=4$ 且 $\Delta_7=\frac{3}{2}$，相应的 7-膜为一个 $III$ 型。
\end{itemize}

% [Tutor Note: 这三条与表 6 的 $\Delta_7$ 列逐一对应：$\Delta_7=3/2$ 是 $IV$ 型（$\SU(3)$）、$2$ 是 $I_0^*$ 型（$\SO(8)$）、$3$ 是 $IV^*$ 型（$E_6$）、$4/3$ 是 $III$ 型（$\SU(2)$）。于是：$k=2$ 可配 $IV,I_0^*,IV^*$；$k=3$ 可配 $III,I_0^*$；$k=4$ 只能配 $III$。注意同一 $k$ 下 $k\Delta_7$ 都落在 $\{2,3,4,6\}$：例如 $k=2,\Delta_7=3/2\Rightarrow k\Delta_7=3$，$k=3,\Delta_7=2\Rightarrow 6$，$k=4,\Delta_7=3/2\Rightarrow 6$。]
正如上面所暗示的，这些 F 理论构造的强大之处在于：人们可以从椭圆几何中读出这些 4d $\cN=2$ SCFT 的各种性质，包括库仑分支几何（秩一情形）、共形反常系数 $(a,c)$、整体对称性，等等。

让我们以整体对称性为例来说明这一点。容易看出，整体对称性有两个明显不同的来源，它们在几何描述中显露出来：一个来自背景的等距，另一个来自驻留在 7-膜上的规范对称性 $F$。垂直于 7-膜的复平面所具有的显然的 $\mathrm{U}(1)$ 等距，被认定为 4d 超共形代数的 $\mathrm{U}(1)_r$ 对称性。7-膜在垂直于 D3 膜的方向上给出 $\C^2/\Z_k$，其等距群在 $k>2$ 时为 $\mathrm{U}(2)$、在 $k=2$ 时为 $\SO(4)$。其中 $\SU(2)$ 子群被认定为 4d 超共形代数中的 $\SU(2)_R$ 对称性，而其余交换子（commutant）则成为味对称性。因此，$k>2$ 时我们有一个 $\mathrm{U}(1)$ 味对称性，$k=2$ 时有一个 $\SU(2)$ 味对称性。至于 7-膜上的规范对称性 $F$，需要考虑 $\Z_k$ 商：只有 $\Z_k$ 不变的子群 $H\subset F$ 才作为 4d SCFT 中的整体对称性存活下来。结合没有 S-折的情形，我们把秩一 4d $\cN=2$ SCFT 的构造汇总在表 8 中。

% [Tutor Note: 这段的对称性账可以逐项验算：$\C^2/\Z_k$（$k>2$）的等距群是 $\mathrm{U}(2)\cong \SU(2)\times \mathrm{U}(1)/\Z_2$，其中 $\SU(2)$ 部分与超共形代数里的 $\SU(2)_R$ 认同，剩下的 $\mathrm{U}(1)$ 成为味对称性；而 $k=2$ 时 $\C^2/\Z_2$ 的等距群是 $\SO(4)\cong \SU(2)\times\SU(2)$，一个 $\SU(2)$ 给 $\SU(2)_R$，另一个 $\SU(2)$ 给味对称性。这与表 8 里 $k=2$ 列出现 $\SU(2)$ 味因子的情形相符。别忘了同时取 $\Z_k$ 不变的子群 $H\subset F$——这一步由表中 $H_0,H_1,H_2$ 等行的上标与 $\rtimes\Z_2$ 记号体现。]

\begin{center}
\small
\begin{tabular}{cccccc}
\toprule
$G$ & $\Delta_7$ & $k=1$ & $k=2$ & $k=3$ & $k=4$ \\
\midrule
$E_8$ & $6$   & $[II^*,E_8]$   & & & \\
$E_7$ & $4$   & $[III^*,E_7]$  & & & \\
$E_6$ & $3$   & $[IV^*,E_6]$   & $[II^*,C_5]$ & & \\
$D_4$ & $2$   & $[I_0^*,D_4]$  & $[III^*,C_3C_1]$ & $[II^*,A_3\rtimes\Z_2]$ & \\
$H_2$ & $3/2$ & $[IV,H_2]$     & $[IV^*,C_2U_1]$  & & $[II^*,A_2\rtimes\Z_2]$ \\
$H_1$ & $4/3$ & $[III,H_1]$    & & $[III^*,A_1U_1\rtimes\Z_2]$ & \\
$H_0$ & $6/5$ & $[II^*,H_0]$   & & & \\
$\varnothing$ & $1$ & & $[I_0^*,C_1\chi_0]$ & $[IV^*,U_1]$ & $[III^*,U_1\rtimes\Z_2]$ \\
\bottomrule
\end{tabular}
\end{center}
\tabtitle{8}{秩一 SCFT（也见表 3），按 [286] 的方式标记，并给出它们在 F 理论中的实现，改自 [57]。$k=1$ 列表示没有 S-折的情形，这些情形已在本节开头介绍过。$\Delta_7=1$ 的最后一行对应于没有 7-膜的情形，因此超对称增强到 $\cN\ge 3$。}

% [Tutor Note: 表 8 的列归属已逐格核对（600 dpi 裁图）：$G$ 与 $\Delta_7$ 两列固定后，每个理论只出现在唯一一个 $k$ 列里——$E_8,E_7,H_0$ 只在 $k=1$；$E_6$ 在 $k=1$（$[IV^*,E_6]$）与 $k=2$（$[II^*,C_5]$）；$D_4$ 在 $k=1,2,3$；$H_2$ 在 $k=1,2,4$（$k=3$ 空）；$H_1$ 在 $k=1,3$（$k=2,4$ 空）；$\varnothing$ 行在 $k=2,3,4$（$k=1$ 空）。原文表中并没有跨行合并的单元格，故此处未用 multirow——若强行合并会破坏「一个理论一格」的对应关系。]
% [Tutor Note: 记号的读法：$[X,Y]$ 中 $X$ 是该理论在 F 理论里所探测的 Kodaira 奇点型（$II^*,III^*,IV^*,I_0^*,IV,III$），$Y$ 是味对称性代数的记号（$A_n$、$C_n$、$U_1$、$\chi_0$ 等，沿用 [286]，与前面的表 3 一致）；$\rtimes\Z_2$ 表示该味代数再与一个 $\Z_2$ 作半直积。$\varnothing$ 一行即「无 7-膜」：$\Delta_7=1$，超对称从 $\cN=2$ 增强到 $\cN\ge3$，与表 7 的 $k=2$ 行（$\tau$ free）呼应。]

类似于 [263,287] 中讨论的 $\cN=3$ S-折，人们可以把上述 $\cN=2$ S-折进一步与离散挠（discrete torsion）结合，也就是在轨道折叠的不动点处加入被困住的 3-形式磁通。在这种情形下，所得 SCFT 的整体对称性与上面的 $\cN=2$ S-折略有不同，但同样可以从椭圆几何中读出。我们这里不涉及这一方面，而建议读者参阅 [59,284]。

% [Tutor Note: 「被困住的 3-形式磁通」指把 $H_3$ 通量的周期取在轨道折叠不动点处的非平凡值（离散挠）。它在 4d 里的效果是改变离散 $\theta$ 角/离散整体对称性的选择，因此「整体对称性略有不同」；这也解释了为什么 [59,284] 能用希格斯化从无挠 S-折得到这些理论。]

\SEC{6}{结论与开放问题}

在正文中，我们基于 2020 年“Quantum Field Theories and Geometry”暑期学校的三次讲座，给出了 4d $\cN=2$ 超共形场论的一个导引。
% [原文 p.93]
正如引言中已经指出的，4d $\cN=2$ SCFT 这一主题在过去三十年间一直是一个活跃而充满活力的研究领域，要把所有有趣的方面都塞进这篇短综述里是不可能的。因此，我们试图为读者提供一本要领手册（guidebook）：它涵盖这些 SCFT 的基本特征，以及在弦/M/F 理论中构造它们的基本工具。我们请读者参阅其他已有的综述以获取互补的视角，以及其中所引文献里更多的细节。最后，尽管在理解这些 SCFT 方面已有实质性进展，仍然存在大量未决问题。下面我们简要讨论其中若干。
% [Tutor Note: 原文 §6「Conclusion and Open Questions」标题实际落在 PDF p.93 底部（印刷 p.92），紧随其后的首段跨页延续到 p.94；BRIEF 里写的「§6 在 p.94 中段」与页图不符（已用页图与文字层双重核对）。另：§1 与 §6 是全书唯一没有编号公式的两节，本节确实不含 (6.x)。]

\para{离散对称性与反常}
到目前为止，在 $\cN=2$ SCFT 中我们只追踪了连续对称性及其 't Hooft 反常。近年来，人们对量子场论中离散对称性及其反常的理解有了很大发展，这些进展在约束 RG 流、刻画红外相图方面显示出强大的威力（例如参见 [259,295]）。因此，搞清这类对称性在 $\cN=2$ SCFT 中如何实现，显然是很重要的；尤其是从弦/M/F 理论的几何构造来看——在那里反常由流入机制（inflow mechanism）[296] 决定（该机制可推广到离散对称性）。近来沿着这条路线已有一些激动人心的进展，涉及从几何构造得到的离散高形式对称性 [278,297--301]。

% [Tutor Note: 「流入机制」指把 $d+1$ 维拓扑作用的边界反常通过 Chern--Simons 型项流入 $d$ 维：反常系数由几何（这里是椭圆纤维化的相交数据）给出。对连续对称性，这给出 't Hooft 反常系数；对离散对称性，同样的机制在 $d+1$ 维的 $\Z_n$ 拓扑项上运作。这正是为什么作者说「几何构造把反常定死」——几何一旦给定，反常就没有自由度了。]

% [原文 p.94]
\para{完全确定一个 4d $\cN=2$ SCFT 所需的数据}
一个 CFT 由它的算子谱、连同满足幺正性、共形不变性与交叉对称性等相容条件的 OPE，来完全确定。然而，当 CFT 具有额外结构（例如 4d $\cN=2$ 超共形对称性）时，这种做法并不高效。在这种情形下，人们可以期望理论由少得多的参数完全确定，例如共形（与味）中心荷以及受保护的局域算子谱，只差一些对缺陷算子谱敏感的离散类。类似的问题已在 [7] 中就规范理论作过分析；理解 $\cN=2$ SCFT 中的这类离散参数（尤其是那些非拉格朗日量的理论）将会很有意义。

\para{4d $\cN=2$ SCFT 的有理性}
4d $\cN=2$ SCFT 的一个奇特特征，是在许多物理可观测量中表现出的有理性程度，而这一点至今没有解释。例如，在已知的 $\cN=2$ SCFT 中，半 BPS 的库仑分支算子全都具有有理的标度维数，与之相关地，共形中心荷也是有理数；这与 $\cN=1$ SCFT 的情形形成对照——在后者中这些量一般是无理数。

% [Tutor Note: 「半 BPS 库仑分支算子的维数有理」与「$a,c$ 有理」这两件事是绑在一起的：$\cN=2$ 的 Weyl 反常满足 Shapere--Tachikawa 型关系，即 $a$、$c$ 可写成库仑分支算子（按 $\Delta_i$）与味对称性流两类贡献之和，故 $\Delta_i$ 的有理性会传导到 $a,c$ 上。$\cN=1$ 没有库仑分支算子的这种保护，$a,c$ 一般是无理数——这正是本段强调的对照。]

\para{极小的 4d $\cN=2$ SCFT}
所有已知的相互作用 4d $\cN=2$ SCFT 的秩都 $\ge 1$，也就是说，至少存在一个半 BPS 的库仑分支算子。然而，先验地并没有理由排除存在一个秩为 0、不含任何半 BPS 库仑分支算子的相互作用 SCFT。具有同等超对称量的这种 SCFT 已知在 3d 中存在，它既没有库仑分支、也没有希格斯分支的半 BPS 算子 [302]。借助超共形自举、通过研究应力张量多重态的四点函数，在 4d 中排除或确认这种可能性，将会很有意义。

% [Tutor Note: 原文此处写 "an interaction SCFT of rank 0"（600 dpi 裁图逐词复核，确为 interaction，缺 -ing；同段其他地方用的是 interacting）。译文按 interacting 处理。物理上要紧的是：秩 0 意味着库仑分支上没有半 BPS 算子（$\cN=2$ 的库仑分支算子都是半 BPS 的），于是「秩 0」与「无半 BPS 库仑分支算子」等价，这正是本段要问的可否存在性。]
\para{库仑分支 EFT 的分类}
库仑分支 EFT 一直是研究 4d $\cN=2$ SCFT 的丰饶园地，因为它为强耦合不动点处的许多物理可观测量提供了一扇窗口。近来已有人尝试通过分析可容许的（admissible）库仑分支 EFT 来分类 $\cN=2$ SCFT [118,145,163,165,166,303]。确实，对于秩一情形，这一纲领成功地给出了一份穷尽的候选理论清单，而这些候选如今都已有紫外完备的构造。然而，其中某些假设与微妙之处仍需澄清。例如，库仑分支手征环被假定为自由生成的（[162,304,305] 中找到了一些并非如此的例子）；此外，离散对称性及其反常尚未被纳入库仑分支 EFT。

% [Tutor Note: 「库仑分支手征环自由生成」这条假设很关键：分类纲领用 Coulomb 分支的算符（$\Delta$ 与 $r$ 荷）作为输入数据，只有在手征环是自由多项式环时，「给定维数谱」才等价于「给定一个环」。一旦手征环有非平凡关系（relation），同一组维数可以对应多个环，分类就会出现歧义——这正是 [162,304,305] 的反例要点。]

% [原文 p.95]
\para{希格斯分支与 3d 镜像拉格朗日量的研究}
对于由 IIB 型几何工程产生的许多强耦合 4d SCFT，希格斯分支较难触及，因为它编码的是三重复形奇点的解消（resolution），而不是复结构形变。正如 4.8.2 节所提到的，研究这类非拉格朗日量理论的希格斯分支的一条途径，是考虑它们的 $S^1$ 紧化，并在镜像框架中寻找所得 3d $\cN=4$ SCFT 的拉格朗日量描述 [199]。对于 $(A,A)$、$(A,D)$ 与 $(D,D)$ 型的 AD 理论类，这件事已在 [234,237,238,240,306] 中被系统地完成。\footnote{[237,238,306] 的作者还研究了这类理论的共形流形；他们还给出了 $D_p(\SU)$ 与 $D_p(\SO)$ 理论的 3d 镜像理论 [43,229]。}继续研究其余 $(G,G')$ 理论的 3d 镜像理论，尤其是那些不承认 Class S 描述的理论类，将会很有意义。

% [Tutor Note: 这里的分工值得记牢：几何工程的复结构形变对应库仑分支，解消（Kähler 形变）对应希格斯分支。所以对三重复形奇点，「读库仑分支」容易（形变参数就是 CB 算子），「读希格斯分支」难（要解消、要处理非孤立/非 crepant 的情形）。$S^1$ 紧化后 3d $\cN=4$ 有镜像描述，希格斯分支在镜像侧变成库仑分支，于是可用磁箭图计算——这就是本段的逻辑链。]

这种构造的一个推广/弱化版本，就是所谓的磁箭图（magnetic quiver）。这些箭图在被理解为 3d 理论时，其库仑分支与我们关心的 4d 理论的希格斯分支重合（同样的概念也适用于 5 维或 6 维理论）[307--312]。\footnote{注意，与 3d 镜像对称不同，我们并不要求磁箭图的希格斯分支与原理论的库仑分支重合。因此，磁箭图构造与其说是两个理论之间的物理联系，不如说是针对几何的一套记账工具。例如，可能出现这样的情况：所关心理论的希格斯分支是若干个超 Kähler 锥的并，而每一个都各自有与之伴随的磁箭图。}它们可以从理论的膜构造 [312--319]\footnote{[319] 还研究了同一希格斯分支的不同但在红外对偶的磁箭图。这一点通过考察超对称线缺陷在对偶下如何变换而得到进一步支持。关于 3d 对偶如何作用于超对称线缺陷的先前工作，包括阿贝尔规范理论中的 [320,321] 与非阿贝尔规范理论中的 [322--325]。}或从它的几何工程 [60,236,264,285,326,327] 计算出来，并可用于获取理论希格斯分支的信息，例如它的对称性、辛奇点的分层（stratification）等等 [318,328--334]。

\para{从几何得到手征代数}
任何 4d $\cN=2$ SCFT 的完整算子代数中都包含一个重要而丰富的子区域（subsector），它由 2d 手征代数 [86] 描述，编码了该 SCFT 的许多特征，尤其是与希格斯分支相关的那些。这就引出一个自然的问题：哪些 2d 手征代数可以出现在 4d $\cN=2$ SCFT 中？既然大多数已知的 SCFT 都来自弦/M/F 理论中的几何构造，人们自然希望直接从几何中辨认出这个手征代数部分。

% [Tutor Note: 2d 手征代数是通过「取 $\cN=2$ 超共形代数里的某个 nilpotent 超荷 $Q$ 的上同调」得到的：$Q$-闭且模去 $Q$-正合的算子构成一个 2d 手征代数，其 $c_{2d}=-12c_{4d}$。它特别能抓住希格斯分支（因为 Higgs 分支的算子都是 $Q$-闭的），而对库仑分支几乎无感——这就是为什么本段开头说它「尤其编码与希格斯分支相关的特征」。]
% [原文 p.96]

\section*{致谢}

作者们感谢 Ibou Bah、Jonathan Heckman、Ken Intriligator、Sara Pasquetti、Shlomo Razamat、Sakura Schafer-Nameki 与 Alessandro Tomasiello 组织了 2020 年暑期学校“QFT and Geometry”，并感谢他们对本文草稿提出的意见。我们还感谢 Fabio Apruzzi、Ibou Bah、Federico Carta、Jacques Distler、Ori Ganor、Simone Giacomelli、Prem Kumar、Noppadol Mekareeya、Carlos Núñez、Diego Rodríguez-Gómez、Lucas Schepers 与 Dan Xie 对草稿的评论和有趣的讨论。最后，我们感谢 Rajeev Singh 在本项目开始阶段的合作。M.A. 部分受 STFC consolidated grant ST/S505778/1 资助。G.A.-T. 受西班牙政府奖学金 MCIU-19-FPU18/02221 资助。A.M. 在 2021 年 9 月之前受“la Caixa”基金会（ID 100010434）资助，奖学金代码 LCF/BQ/IN18/11660045，并受欧盟“地平线 2020”研究与创新计划 Marie Skłodowska-Curie 资助协议 No.\ 713673 资助。A.M. 的工作还部分得到德国研究基金会（DFG）在德国卓越战略 EXC 2121 “Quantum Universe”390833306 之下的支持。H.-Y.S. 部分受西蒙斯基金会 Simons Collaborations on Ultra-Quantum Matter（grant No.\ 651440 (A.K.)）资助。Z.S. 感谢美国能源部（DOE）在合作协议 DE-SC0009919 之下的支持，以及西蒙斯基金会 award \#568420 的支持。Y.W. 部分受哈佛大学数学科学与应用中心（Center for Mathematical Sciences and Applications）与自然基本规律中心（Center for the Fundamental Laws of Nature）资助。

% [边界] 本块到此结束：原文 p.97（印刷 p.96）上部「致谢」段末。自下一行起原文为 References，
% 由 Lead 的脚本生成（见 BRIEF §2 表注），tr-08 未撰写任何文献条目。
% [Tutor Note: 致谢中的基金编号与人名一律保留原文（专有名词不译）；此处逐字核对过 500+ dpi 页图，注意 "Marie Skłodowska-Curie" 的 ł、"Núñez"/"Rodríguez-Gómez" 的重音符号，以及 grant 编号 ST/S505778/1、MCIU-19-FPU18/02221、LCF/BQ/IN18/11660045、713673、EXC 2121、390833306、651440、DE-SC0009919、\#568420 均照原文。]
